% Options for packages loaded elsewhere
% Options for packages loaded elsewhere
\PassOptionsToPackage{unicode}{hyperref}
\PassOptionsToPackage{hyphens}{url}
\PassOptionsToPackage{dvipsnames,svgnames,x11names}{xcolor}
%
\documentclass[
  10pt,
  letterpaper,
  nols]{tufte-book}
\usepackage{xcolor}
\usepackage{amsmath,amssymb}
\setcounter{secnumdepth}{5}
\usepackage{iftex}
\ifPDFTeX
  \usepackage[T1]{fontenc}
  \usepackage[utf8]{inputenc}
  \usepackage{textcomp} % provide euro and other symbols
\else % if luatex or xetex
  \usepackage{unicode-math} % this also loads fontspec
  \defaultfontfeatures{Scale=MatchLowercase}
  \defaultfontfeatures[\rmfamily]{Ligatures=TeX,Scale=1}
\fi
\usepackage{lmodern}
\ifPDFTeX\else
  % xetex/luatex font selection
  \setmainfont[]{Times New Roman}
  \setsansfont[]{Times New Roman}
\fi
% Use upquote if available, for straight quotes in verbatim environments
\IfFileExists{upquote.sty}{\usepackage{upquote}}{}
\IfFileExists{microtype.sty}{% use microtype if available
  \usepackage[]{microtype}
  \UseMicrotypeSet[protrusion]{basicmath} % disable protrusion for tt fonts
}{}
\makeatletter
\@ifundefined{KOMAClassName}{% if non-KOMA class
  \IfFileExists{parskip.sty}{%
    \usepackage{parskip}
  }{% else
    \setlength{\parindent}{0pt}
    \setlength{\parskip}{6pt plus 2pt minus 1pt}}
}{% if KOMA class
  \KOMAoptions{parskip=half}}
\makeatother
% Make \paragraph and \subparagraph free-standing
\makeatletter
\ifx\paragraph\undefined\else
  \let\oldparagraph\paragraph
  \renewcommand{\paragraph}{
    \@ifstar
      \xxxParagraphStar
      \xxxParagraphNoStar
  }
  \newcommand{\xxxParagraphStar}[1]{\oldparagraph*{#1}\mbox{}}
  \newcommand{\xxxParagraphNoStar}[1]{\oldparagraph{#1}\mbox{}}
\fi
\ifx\subparagraph\undefined\else
  \let\oldsubparagraph\subparagraph
  \renewcommand{\subparagraph}{
    \@ifstar
      \xxxSubParagraphStar
      \xxxSubParagraphNoStar
  }
  \newcommand{\xxxSubParagraphStar}[1]{\oldsubparagraph*{#1}\mbox{}}
  \newcommand{\xxxSubParagraphNoStar}[1]{\oldsubparagraph{#1}\mbox{}}
\fi
\makeatother


\usepackage{longtable,booktabs,array}
\usepackage{calc} % for calculating minipage widths
% Correct order of tables after \paragraph or \subparagraph
\usepackage{etoolbox}
\makeatletter
\patchcmd\longtable{\par}{\if@noskipsec\mbox{}\fi\par}{}{}
\makeatother
% Allow footnotes in longtable head/foot
\IfFileExists{footnotehyper.sty}{\usepackage{footnotehyper}}{\usepackage{footnote}}
\makesavenoteenv{longtable}
\usepackage{graphicx}
\makeatletter
\newsavebox\pandoc@box
\newcommand*\pandocbounded[1]{% scales image to fit in text height/width
  \sbox\pandoc@box{#1}%
  \Gscale@div\@tempa{\textheight}{\dimexpr\ht\pandoc@box+\dp\pandoc@box\relax}%
  \Gscale@div\@tempb{\linewidth}{\wd\pandoc@box}%
  \ifdim\@tempb\p@<\@tempa\p@\let\@tempa\@tempb\fi% select the smaller of both
  \ifdim\@tempa\p@<\p@\scalebox{\@tempa}{\usebox\pandoc@box}%
  \else\usebox{\pandoc@box}%
  \fi%
}
% Set default figure placement to htbp
\def\fps@figure{htbp}
\makeatother





\setlength{\emergencystretch}{3em} % prevent overfull lines

\providecommand{\tightlist}{%
  \setlength{\itemsep}{0pt}\setlength{\parskip}{0pt}}



 


\usepackage{booktabs}
\usepackage{longtable}
\usepackage{array}
\usepackage{multirow}
\usepackage{wrapfig}
\usepackage{float}
\usepackage{colortbl}
\usepackage{pdflscape}
\usepackage{tabularx}
\usepackage{xltabular}
\usepackage{threeparttable}
\usepackage{threeparttablex}
\usepackage[normalem]{ulem}
\usepackage{makecell}
\usepackage{xcolor}
\usepackage{fvextra}
\DefineVerbatimEnvironment{Highlighting}{Verbatim}{
  breaklines,
  breakanywhere=true,
  commandchars=\\\{\}
}
\makeatletter
\@ifpackageloaded{bookmark}{}{\usepackage{bookmark}}
\makeatother
\makeatletter
\@ifpackageloaded{caption}{}{\usepackage{caption}}
\AtBeginDocument{%
\ifdefined\contentsname
  \renewcommand*\contentsname{Table of contents}
\else
  \newcommand\contentsname{Table of contents}
\fi
\ifdefined\listfigurename
  \renewcommand*\listfigurename{List of Figures}
\else
  \newcommand\listfigurename{List of Figures}
\fi
\ifdefined\listtablename
  \renewcommand*\listtablename{List of Tables}
\else
  \newcommand\listtablename{List of Tables}
\fi
\ifdefined\figurename
  \renewcommand*\figurename{Figure}
\else
  \newcommand\figurename{Figure}
\fi
\ifdefined\tablename
  \renewcommand*\tablename{Table}
\else
  \newcommand\tablename{Table}
\fi
}
\@ifpackageloaded{float}{}{\usepackage{float}}
\floatstyle{ruled}
\@ifundefined{c@chapter}{\newfloat{codelisting}{h}{lop}}{\newfloat{codelisting}{h}{lop}[chapter]}
\floatname{codelisting}{Listing}
\newcommand*\listoflistings{\listof{codelisting}{List of Listings}}
\makeatother
\makeatletter
\makeatother
\makeatletter
\@ifpackageloaded{caption}{}{\usepackage{caption}}
\@ifpackageloaded{subcaption}{}{\usepackage{subcaption}}
\makeatother
\usepackage{bookmark}
\IfFileExists{xurl.sty}{\usepackage{xurl}}{} % add URL line breaks if available
\urlstyle{same}
\hypersetup{
  pdftitle={Primer of Ecology},
  pdfauthor={Hank Stevens},
  colorlinks=true,
  linkcolor={blue},
  filecolor={Maroon},
  citecolor={Blue},
  urlcolor={Blue},
  pdfcreator={LaTeX via pandoc}}


\title{Primer of Ecology}
\author{Hank Stevens}
\date{2026-08-19}
\begin{document}
\frontmatter
\maketitle

\renewcommand*\contentsname{Table of contents}
{
\hypersetup{linkcolor=}
\setcounter{tocdepth}{2}
\tableofcontents
}

\mainmatter
\bookmarksetup{startatroot}

\chapter{Primer of Ecology using R}\label{primer-of-ecology-using-r}

\bookmarksetup{startatroot}

\chapter*{Preface}\label{preface}
\addcontentsline{toc}{chapter}{Preface}

\markboth{Preface}{Preface}

UPDATED 2023-08-10

Chapter numbers have changed, starting after Lotka-Volterra competition.

Updated original, but long absent, material on the
competition-colonization tradeoff, extinction debt (sensu Tilman 1994),
and especially the \textbf{storage effect}.

Later this fall, I may add a section on \emph{modern coexistence
theory}.

August 2022

These updates will include usual fixes for typos and incremental
movement toward slightly practical application. I am trying to include
slightly more material and code that students might actually use in
their theses and dissertations. When I update the book, I will list
below as many of the changes as I can, including

\begin{itemize}
\tightlist
\item
  cleaning up and clarifying the integral projection example.
\item
  includes a simple discrete nonlinear density-dependent population
  growth model (Watkinson et al.~1980). Several studies evaluating
  modern coexistence theory and the invasion criterion of coexistence
  use extensions of this model.
\item
  I will soon add nonlinear density-dependent population growth model
  (Watkinson et al.~1980) into two species competition.
\end{itemize}

20 August 2021

NEW IN THIS EDITION OF THE E-BOOK

\begin{itemize}
\tightlist
\item
  Lotka-Volterra versions of direct competition AND mutualism are
  together in one chapter, followed by consumer-resource interactions,
  including predator-prey, and resource-based competition and mutualism.
\item
  a chapter on consumer-resource competition and mutualism that follows
  chapters covering two-species consumer resource interactions (mostly
  predator-prey).
\item
  If it has been \textgreater{} 1 y since looking at the \texttt{primer}
  package, you might want to peruse the index - I added a few more
  functions including ones for 1/f noise, spectral mimicry, and the
  Allee effect.
\end{itemize}

\begin{longtable}[]{@{}
  >{\raggedright\arraybackslash}p{(\linewidth - 0\tabcolsep) * \real{0.0556}}@{}}
\toprule\noalign{}
\begin{minipage}[b]{\linewidth}\raggedright
This is based on my book published with Springer, ``A Primer of Ecology
With R'', part of the Use R! series.
\end{minipage} \\
\midrule\noalign{}
\endhead
\bottomrule\noalign{}
\endlastfoot
title: Theory in Ecology \\
\end{longtable}

\begin{verbatim}




::: {.cell}

```{.r .cell-code}
library(tidyverse)
library(dagitty)
library(primer)
```
:::


In this chapter, we introduce a perspective on ecological theory, and
provide two examples of efficient theory, metabolic scaling and maximum
entropy theory.

Theory is really, really important (@fig-theorydatanathist). Without theory of some sort to guide
us, we are merely counting pebbles in a quarry. Good theory may be
conceptual or mathematical, and helps us focus more clearly and in a
less ambiguous manner on the systems we study. In this book, we explore
mathematical models that provide special cases of theory.


::: {.cell .column-margin}

```{.r .cell-code}
g <- dagitty(
  "dag{
  Theory -> Natural_History -> Data;
  Data -> Theory;
  Theory <- Natural_History <- Data;
  Data <- Theory
  }"
)
coordinates(g) <- list(
  x=c(Theory=1, Natural_History=2, Data = 0),
  y=c(Theory=0, Natural_History=sqrt(3), Data = sqrt(3))
)
plot(g)
```

::: {.cell-output-display}
![Ecology and evolutionary biology require all of these to advance.](01-theory_files/figure-pdf/fig-theorydatanathist-1.pdf){#fig-theorydatanathist fig-pos='H'}
:::
:::


*Scientific theory* is a body of knowledge that provides an organized
and mechanistic view of how the world works [@scheiner2010]. Theories
concerning gravity, general relativity, and evolution by natural
selection provide structured ways of connecting observations, patterns,
and processes that provide insight into why the world is the way it is.
This stands in stark contrast to the colloquial use of *theory* that
implies a lack of knowledge, as when someone says "oh, that's just a
theory", referring to a guess without much evidence. Scientific theory
is a set of explanations whose validity has been tested repeatedly by
experiments and new data.

## Examples of theories

Ecology has lots of theories, of all different types. Below I discuss
some which may be prevalent, important, useful, or some combination.

### Hierarchy theory

![A hierarchical perspective called Macrosystems Theory (Rose et al.
(2017))](figs/Rose_et_al_2017_fig_1.pdf){width="80%"}

An early and persistent organizing schema in ecology is based on
*hierarchy theory* [@oneill1986; @rose2017, and references therein]. It
posits that ecological systems are structured *hierachically*, such that
each entity comprises subunits. For instance, an entity such as a
population of big bluestem grass (*Andropogon gerardii*) is part of a
larger ecology community of many species. The population of big bluestem
comprises subpopulations separated in space, a subpopulation comprises
separate individuals, that each individual comprises multiple ramets and
a set of organ systems and tissues, which comprise different cell types.
This theory posits that each entity gives rise to emergent properties to
the hierarchical level above it, and influences processes within each
smaller sub-entity in the hierachical level below it. As a disciplinary
organizing principle, this approach structures nearly all of the ecology
curriculum.

Hierarchy theory gets more complicated when the levels of a hierarchy
start to include fundamentally different types of entities. The big
bluestem hierarchy above included only biotic components--a individual
is part of a population which is part of a community of individuals of
multiple species, and is made up of organ systems and tissues. Ecology,
however, includes both the biotic and the abiotic parts of environments.
An ecosystem includes a community of species, but also the nutrients,
water, light, and other abiotic components, along with the spatial
arrangement of all of these things.

Different hierarchies are useful for different questions. An individual
organism can play very different roles in different hierarchies.
Consider and individual bunch grass. To understand how a population
evolves, we need to count individuals within a population, because
evolutionary fitness is tracked by the number of independent
reproductive units. In contrast, to understand competitive interactions,
it may be much more important to weigh the biomass of groups of
individuals in a population, because biomass is more closely related of
resource uptake.

### A general theory of ecology

Good scientific theories exist within a hierarchy of disciplinary
knowledge [@scheiner2011]. They explain phenomena within a *domain* of
knowledge which is organized around *principles* and *assumptions*.
Scheiner and Willig posit a theory of biology that explains phenomena
relating to the "diversity and complexity of living systems". One of the
ten principles on which this theory depends is that "the cell is the
fundamental unit of life". Subsumed within their theory of biology is
the theory of cells whose domain is "cells and the causes of their
structure, function, and variation." This theory in turn is based on
principles and has theories to organize our understanding of cells and
what cells do.

*Models* are specific and explicit manifestations of more general
theories. In this book, we focus on popular mathematical models that are
specific manifestions of theories of ecology.

@scheiner2011 propose a theory of ecology, some of which we cover in
this book. Here is part of this theory:

**The General Theory of Ecology**

**Domain:** The spatial and temporal pattern of the distribution and
abundance of organisms, including causes and consequences.

**Principles:**

1.  Organisms are distributed in space and time in a heterogeneous
    manner.
2.  Organisms interact with their abiotic and biotic environments.
3.  Variation in the characteristics of organisms results in
    heterogeneity of ecological patterns and processes.
4.  The distributions of organisms and their interactions depend on
    contingencies.
5.  Environmental conditions as perceived by organisms are heterogeneous
    in space and time.
6.  Resources as perceived by organisms are finite and heterogeneous in
    space and time.
7.  Birth and death rates are a consequence of interactions with the
    abiotic and biotic environment.
8.  The ecological properties of species are the result of evolution.

These principles constitute what we know is true about ecological
systems. Some of these principles provide the focus for a single chapter
while other principles apply broadly to many chapters in this book.

Here is my own perspective on a general theory of ecology:

*Domain:* The house of life[^01-theory-1]: its constituent entities,
causes, and consequences.

[^01-theory-1]: "ecology" derives from the Greek "oikos" which means the
    rules of the house

*Principles:*

1.  Entities[^01-theory-2] are open systems with inputs and outputs.
2.  Entities have internal complexity.
3.  Entities include self-replicating components (living elements).
4.  Entities interact via inputs, outputs, and behavior.
5.  Rates of change, including inputs and outputs, are influenced
    directly by physical factors: space, temperature, and concentration.

[^01-theory-2]: An entity may be an ecosystem, a community, a
    population, an individual, or some other system with operationally
    defined boundaries.

You will see elements of these principles throughout this book as well.

### Efficient theory

@marquet2014 argue that the best theories are those which are
*efficient*. Such theories tend to be *based on first principles*, which
are observations and laws that are fundamental assumptions in a
scientific domain. In biology, such principles can include the laws of
thermodynamics, and mathematical properties such as the central limit
theorem. Theories built upon first principles are thus well-grounded in
reality as we understand it and lead logically to refinements. Marquet
and his colleagues also claim that efficient theory is *expressed in
mathematics*. Mathematics is a universal language that is unamibiguous.
It forces us to be as clear as possible about what we mean when we state
a theory.[^01-theory-3] Last, efficient theories are those that make a
*large number of predictions using only a small number of free
parameters*.[^01-theory-4] Examples of efficient theories we cover in
this book include metabolic scaling, exponential growth, density
dependence, and ecological neutral theory.

[^01-theory-3]: $E=mc^2$ - need I say more?

[^01-theory-4]: Variables are quantities we measure and which change
    through time (e.g., population size). *Parameters* are (usually)
    fixed constants that govern the rates of change of variables (e.g.,
    per capita birth rate).

Marquet et al. and Scheiner and Willig emphasize slightly different
features of the definition of "theory". Scheiner and Willig emphasize
relatively broad ideas that are well-supported by experiments and
repeated observation. Marquet and colleagues tend to mean something
fairly specific and narrow, typically something that can be expressed
mathematically. Scheiner and Willig might refer to such theory as
constitutive theory or even simply a model.

Next, I describe the Metabolic Theory of Ecology. This theory is based
on first principles, and its central tenets are expressed
mathematically. It's core equation has a very small number of free
parameters (fitted constants) and makes a very large number of testable
predictions. Parts of this theory are supported by a very large number
of observations. It fits everyone's definition of theory.

## An example: Metabolic Theory of Ecology {#metabolic_scaling}

Metabolic rate is central to how rapidly individuals forage for, consume
and use resources, reproduce and die. The *metabolic theory of ecology*
[@brown2004a] is a well-supported body of knowledge about the underlying
mechanisms, and the resulting profound and wide-ranging consequences for
populations and ecosystems.

Body size and temperature are fundamental properties of organisms and
the environment. The study of how body shape and body processes scale
with body size is *allometry*. Because body size affects metabolic rate,
body size indirectly helps determine population growth rates and how
species interact with each other. Temperature affects how molecules
vigorously molecules vibrate and move, and so increasing temperature
tends to speed up chemical reactions. As metabolism is really just a
complex network of biochemical reactions, temperature influences
metabolic rate.

The core of this theory is expressed in a simple mathematical equation
that describes how body size and temperature govern metabolic rate.

### Body-size dependence

There is a profoundly simple and general rule describing the effect of
interspecific variation in body size on metabolism.[^01-theory-5] This
biological law is referred to as the Kleiber law [@kleiber1932], or
*quarter power scaling* [@brown2004a]. When we compare the basal (i.e.
resting) metabolic rates of different species, across a wide range of
body sizes spanning many orders of magnitude, we find that

[^01-theory-5]: Metabolic rate may be measured by variables tied
    directly to metabolism, such as the rate of oxygen or energy
    consumption, or CO$_2$ production.

*whole-organism resting metabolic rate increases with organism mass
raised to the three-quarter power*, or,

$$
B = aM^{z} \quad;\quad z = 3/4
$$

In this equation, $B$ is basal, or resting, metabolic rate, $M$ is body
mass, $a$ is a proportionality constant, and $z$ is the *power law
scaling coefficient*. The proportionality constant $a$ varies depending
on the type of organism such as arthropods, fish, or mammals. Plants
scale in the same manner [@niklas2001], although size or mass is a
little trickier to measure. The scaling coefficient, $z$, is the
seemingly magical constant that many have argued does not vary
substantially among different types of organisms.

Ecologists typically describe metabolism-mass relations and other power
law behavior using logarithmic scales. When we do that, power law
relations become linear. Using our rules for exponents and logarithms,
metabolic scaling becomes $$ \log B = \log a + z\log M$$ so that
$\log B$ increases linearly with $\log M$ with a slope of $3/4$. Our
brains can process and compare linear relations much more easily than
curvilinear ones.

Here we plot the curvilinear relation in R using `curve()` in the
`graphics` package of R that is included in the base installation as one
of the core packages. The function `curve()` can plot any curve that be
expressed as a function of `x`. Below, we draw a curve of a dotted 1:1
line for comparison, and then *add* the power function $x^{3/4}$.


::: {.cell .column-margin}
::: {.cell-output-display}
![Metabolic rate increases predictably with species body sizes.](01-theory_files/figure-pdf/fig-MTE34-1.pdf){#fig-MTE34}
:::
:::



::: {.cell}

```{.r .cell-code}
## using curve, let your variable be 'x'.
curve(1*x, from = .01, to=100, ylab = "Metabolic rate (B)",
      xlab="Body mass (M)", lty=3)
curve(x^(3/4), from = .01, to = 100, add=TRUE)
```
:::


To help us grasp the implications of this, let's consider
*mass-specific* metabolic rates. "Mass-specific" means on a per-gram
basis.[^01-theory-6] Mass-specific metabolic rate is basal metabolic
rate of an individual divided by its mass, or $B/M$.

[^01-theory-6]: With plants, we often measure something called "specific
    leaf area", SLA, which is the two dimensional area of a leaf divided
    by its mass.

The mass-specific metabolic rate allows us to compare directly, for
example, the metabolic rate of a cell in a shrew vs. a cell in an
elephant. Which cell is burning fuel faster?

We can estimate this from the above metabolic scaling principle and the
using rules exponents
$$ \frac{B}{M} = a \frac{M^z}{M^1} = a M^{z-1} = aM^{-1/4}$$ From this,
we now have the rule that

*mass-specific metabolic rate declines with organisms mass raised to the
negative one quater power*


::: {.cell .column-margin}
::: {.cell-output-display}
![Mass-specific metabolic rate decreases predictably with species body sizes.](01-theory_files/figure-pdf/fig-MTE-MS-1.pdf){#fig-MTE-MS}
:::
:::


Over the years, there has been heated debate about (i) the precise value
of the scaling coefficient $z$, and (ii) the underlying mechanism. Early
arguments suggested that $z \approx 2/3$ because the rate heat
dissipation scales with the amount surface area. Why $2/3$?

Let's envision the volume of an organism having three linear dimensions,
so the volume scales to the cube of linear dimensions, while the surface
area scales to the square of these linear dimensions,[^01-theory-7]
$$V \propto L^3$$ $$A \propto L^2$$ The early explanation was that
metabolic rate, $B$, scales linearly with area,
$$B \propto A^1 \propto L^2$$. With substitution we get,
$$L^2 \propto V^z \propto (L^3)^z$$ implying that the exponents $2 = 3z$
or $z=2/3$, so we get, $$B=V^{2/3}$$, and, for the most part, mass
scales linearly with volume for mammals or any other such group.

[^01-theory-7]: A linear relation is one in which $y$ is proportional to
    $x$, or $y \propto x$.

This early theory was because it started with first principles (heat
dissipation and geometry) and resulted in the prediction of a single
parameter. It could then be used to make predictions about how metabolic
rate scales with body mass. Metabolic rate governs a huge amount of
biology and ecology, including resource consumption rates, lifespan, and
maximum population growth rates. Therefore, this theory and this model
could be powerful tools for understanding the world and making testable
predictions.

The above model is good because it could be tested. That is what has
been done, and scientists found that there was a consistent mismatch
between observations and the theory. Investigators showed that the value
of the exponent appeared closer to 3/4 raher than 2/3. In the 1990s, a
group including Jim Brown and Geoffrey West [@west1997] proposed an
underlying mechanism that explained why it should be 3/4. They assumed
that organisms must

-   distribute resources from a single source through a branching,
    fractal-like, space-filling network to all parts of the body,
-   the size of the smallest branch ( a capillary) was the same for
    organisms of all sizes.
-   the energy required to distribute the resources must be minimized,
    that less energy-efficient designs would be lost through natural
    selection.

The prediction that resulted from these assumptions was that the
exponent would be 3/4. This theory and model begin with different first
principles and makes a different prediction.

Soon Jayanth Banavar and his colleagues offered an alternative
[@banavar1999; @banavar2002], arguing that the assumption of the
fractal-like network was not correct, and in any event, was not
necessary and did not apply to all organisms. They proposed different
theory with less restrictive assumptions and found nonetheless that the
exponent was also predicted to be 3/4.

At the base of all these arguments is the geometry of the resource
distribution system. All organisms take in limiting resources and have
to distribute those resources to each part of each cell in the body. The
key point is that *the larger the organism, the greater the portion of
the resources are in transit at any instant in time*. This leads to an
increasingly inefficient system, in which the metabolism of larger
organisms has to run more slowly per unit resource:

*Larger organisms can process more resources per unit time
(*$B=aM^{3/4}$), but do so less and less efficiently
($\frac{B}{M}=aM^{-1/4}$) due to resources in transport.

Banavar, Brown and others eventually collaborated to address quarter
power scaling in animals in particular which led to additional novel
predictions [@banavar2010].

This theory remains a fertile and active area of research
[@glazier2018]. The interested reader should be careful to distinguish
between patterns observed across many species of very different sizes,
versus patterns observed in a single species with individuals of
different sizes versus other types of patterns. Subtly different
patterns may be driven be very different mechanisms.

### Temperature dependence

In addition to body size, temperature plays the other key role in
regulating metabolic rate. The Arrhenius equation connects the
macroscopic property of temperature to the kinetic energy of molecules
and the rates they govern. Metabolic rate is proportional to these rate
determining processes, $$B  = a e^{\frac{-E_a}{kT}}$$ where $a$ is just
a constant, $e$ is the exponential, $E_a$ is the average activation
energy of rate-limiting enzymes (units, eV), $k$ is Boltzmann's constant
(units eV$\,$K$^{-1}$), and $T$ (units deg K). Bolztmann's constant
($\backsim 8.6 \times 10^-5$$\,$eV$\,$K$^{-1}$) converts the macroscopic
property of temperature to kinetic energy of molecules.

Individual biochemical reactions combine to determine basal metabolic
rate, so @gillooly2000 have taken this as a foundation for the metabolic
theory of ecology [@brown2004a]. In 2000, they suggested that the
average activation energy is approximately $E_a = 0.23\,$eV . Subsequent
work has described this as "temperature sensitivity", where larger
numbers imply that organisms respond more strongly to temperature
variation.

The Arrhenius equation is a little more complicated that a simple power
law, but not too much. Over the range of biologically relevant
temperatures, it is dominated by a largely exponential increase in
metabolic rate with increasing temperature (Fig \@ref(fig:arrh)).


::: {.cell .column-margin}

```{.r .cell-code}
# with base R
# base R: curve(10^4*exp(-0.23/(8.5 * 10^-5 *x)), 276, 316), ylab="B", xlab='T')
# or ggplot2
# the function, with default parameter values
eq.t <- function(t,a=10^4,E=0.23,k=8.6 * 10^-5){a*exp(-E/(k*t))}
# the data used in our function
temps <- data.frame(t=276:316)
ggplot(data=temps, aes(x=t)) + 
# set the basic properties of the plot
  stat_function(fun=eq.t, geom="line") + 
# set the function to plot
  xlab("Temperature (K)") + ylab("Metabolic rate (B)")  
# add labels
```

::: {.cell-output-display}
![The effect of body temperature on ectothermic metabolic rates can be approximated with the Arrhenius function, $B = a e^{-E_a/(kT)}$. Here $a = 10^4$, and $E_a = 0.23$. It is similar in shape to a power law with $z > 1$, over the range of biologically relevant temperatures.](01-theory_files/figure-pdf/fig-arrh-1.pdf){#fig-arrh fig-pos='H'}
:::
:::


When we linearize the relation between metabolic rate and temperature,
we get $$
\begin{aligned}
B &= a e^{\frac{-E_a}{kT}}\\
\log(B) &= \log{a} - E_a\frac{1}{kT}\\
\end{aligned}
$$ where the dependent variable is $1/(kT)$, $-E_a$ is the slope, and
$\log a$ is the intercept. Thus, the negative slope of this relation
describes theoretical prediction for the effect of temperature on
metabolic rate.

**So, there you have it.** The metabolic theory ecology is the algebraic
product of body size- and temperature-dependence:
$$B = a M^{3/4} e^{\frac{-E_a}{kT}}$$ This theory makes quantitative
predictions regarding all kinds of ecology phenomena [@brown2004a],
including

-   home range size
-   population growth
-   population size
-   resource uptake
-   predation and other species interactions, and
-   ecosystem cycling.

Note that these relations are based on first principles of geometry and
thermodynamics, and that they depend on only a small number of
parameters ($a$, $-E_a$, and perhaps $z=3/4$), and makes a tremendous
number of predictions. Therefore, @marquet2014 propose that this is
"good" theory, and very *efficient*.

## Power law scaling implies constant relative differences

In power law scaling, relative change is constant. That is, a
proportional change in one variable results in a proportional change in
the other. For instance, when we compare a smaller species to a larger
species with $100 \times$ the body mass, we can expect to see metabolic
rate increase by $31.6 \times$, *regardless of the mass of the smaller
species.* For now, we will verify this numerically for some limited
cases.


::: {.cell}

```{.r .cell-code}
# define body mass and metabolic rate
m <- c(.01, 1, 100, 10000)
b <- m^.75
```
:::


Now we will divide each mass $i$ by the next smallest mass $i-1$. We do
that using a vector by dividing each mass except the first one, by each
mass except the last one.


::: {.cell}

```{.r .cell-code}
# round(x, digits=0) rounds number to zero decimal places
round( m[-1]/m[-length(m)], digits = 0)
round( b[-1]/b[-length(b)], digits = 1)
```
:::


When we do these divisions, we see the constant relative change (Table
\@ref(tab:relativemb)).


::: {.cell}
::: {.cell-output-display}


Table: As we increase mass by a constant multiplier (10x), power law scaling results in a constant proportional change in basal metabolc rate.

|                     | Small|   Med.|    Big|     Huge|
|:--------------------|-----:|------:|------:|--------:|
|Mass                 |  0.01|   1.00| 100.00| 10000.00|
|Basal.metabolic.rate |  0.03|   1.00|  31.62|  1000.00|
|Relative.change.m    |    NA| 100.00| 100.00|   100.00|
|relative.change.b    |    NA|  31.62|  31.62|    31.62|


:::
:::


We can verify this generally using algebra, not just in the particular
case above. We will show that if mass increases by a constant
multiplier, metabolic rate will also, regardless of the particular
masses involved.

Let mass $m_2$ be greater than mass $m_1$ by a factor of $c$, so that
$m_2 = c m_1$, and $$\frac{m_2}{m_1} = c$$.

We would like to show that the ratio of the metabolic rates $b_2 / b_1$
is also a constant. Since $m_2 = cm_1$, we can say that
$$b_1 = a m_1^{3/4}$$ $$b_2 = a (cm_1)^{3/4} = ac^{3/4}m_1^{3/4}$$
$$\frac{b_2}{b_1} = \frac{ac^{3/4}m_1^{3/4}}{am_1^{3/4}}$$ When we
reduce this fraction, we a left with $$\frac{b_2}{b_1} = c^{3/4}$$

This shows that with power law scaling, increasing $x$ by a constant
*multipier* (or proportion), $y$ increases by the same proportion raised
to that power.

Let's represent this graphically in a couple of ways, reusing data we
made up previously in this chapter. First, we just change the axes
themselves, so that the units of the scales are multiples of 10 (often
in scientific notation).


::: {.cell .column-margin}

```{.r .cell-code}
# using base R
par(mar=c(5,4,0,0), mgp=c(1.5,.4,0) )# set figure margins in "lines"
curve(x^(3/4), from = .01, to = 100, log="xy", ylab="Basal metabolic rate", xlab="Mass")
text(10, 80^.7, expression(M^0.75))
```

::: {.cell-output-display}
![Changing the scales of the axes to linearize power law relations. Note scales are logarithmic, using the original linear values.](01-theory_files/figure-pdf/fig-logscale-1.pdf){#fig-logscale fig-pos='H'}
:::
:::


## Meet METE: Maximum Entropy Theory of Ecology

The Maximum Entropy Theory of Ecology is best known for making
predictions and elucidating patterns of species-area relations and
species abundance distributions that emerge in relatively stable
ecological communities and ecosystems [@harte2014]. However, it is an
active area of research and Harte and others are extending the ideas
into non-equilibrium systems as well [@harte2021mete].

Entropy is, basically, disorder. Less entropy means more order or
structure. Maximum entropy is maximum disorder. The universe tends
toward maximum entropy, although in some places and times, like parts of
our world right now, entropy can decrease. Physicists are the folks most
responsible for our understanding of entropy, and Edwin Jaynes is
arguably the one most responsible our understanding and appreciation of
its role in information and probability.

MaxEnt, or the principle of maximization of information entropy, or the
principle of maximum entropy is that, given a set of constraints on a
system, the least-biased probability distribution that describes data
arising from that system will be the one with maximum entropy.

The principle of maximum entropy has show that maximal uncertainty
arising from the flattest distribution, given known constraints,
provides the least biased description of a system. It can provide the
least biased probability distributions for these situations with these
constraints:

-   Consider a jar contains jelly beans, and contains at least one and
    no more than 1000. What is the least biased distribution of possible
    counts of jelly beans?
-   Imagine a hectare of land that contains trees of different heights,
    and that the mean height is three meters. What is the least biased
    distribution of possible tree heights?
-   Regarding this forest - what if you were told that the variance in
    height was two meters? What would be the least biased distribution
    of possible tree heights?

The principle of maximum entropy could help us prove that, in the
absence of other information, the least biased distributions for these
sets of constraints would be the uniform, exponential, and Gaussian (or
normal) probability distributions.

The maximum entropy theory of ecology [@harte2014] can tell us what to
expect when we know very little, or *when very little is going on aside
from random noise.*

We will talk more about the METE in our chapter on diversity, but I
mention it here to emphasize that patterns arise in ecology via random
noise.

The rank-abundance distribution is a famous ecological pattern. It is
the log-abundance or relative abundance of species, arranged from the
most abundant to the least abundant (Fig. \@ref(fig:mete1)).


::: {.cell .column-margin}

```{.r .cell-code}
library(vegan)
data(BCI)
abun <- colSums(BCI)
mod <- rad.lognormal(abun)
plot(mod, main="Rank-abundance distribution")
```

::: {.cell-output-display}
![Rank abundance distribution of tree species in a 50 ha plot on Barro Colorado Island. Line is a best fit of the log-normal distribution.](01-theory_files/figure-pdf/fig-mete1a-1.pdf){#fig-mete1a fig-pos='H'}
:::
:::


The curve formed by these data points (Fig. @fig-mete1a) is
steep at first, among the most abundant species. The curve then flattens
out a bit before getting a little bit steeper again. 


::: {.cell .column-margin}

```{.r .cell-code}
mod.oct <- prestonfit(colSums(BCI))
plot(mod.oct, main="Species abundance distribution")  
```

::: {.cell-output-display}
![Species abundance distribution of tree species in a 50 ha plot on Barro Colorado Island. Line is a best fit of the log-normal distribution.](01-theory_files/figure-pdf/fig-mete1b-1.pdf){#fig-mete1b fig-pos='H'}
:::
:::


If we group species into logarithmic abundance categories, and make a histogram of
abundances, we get the \*species abundance distribution (Fig. \@ref(fig:mete1b)). @preston1962 noted that the distributions were remarkabling similar to a normal
(Gaussian) distribution, when abundances were log-transformed (resulting
in a lognormal-like distribution). It turns out that this curve is quite
common and for decades ecologists argued about how we might infer the
underlying ecological processes on the basis of the shape of the curve.
This discussion died out once ecologists admitted that it was impossible
to determined the processes on the basis of the pattern.

The discussion of the shape of the rank-abundance distribution was
revived in the late nineties when Steve Hubbell showed that if you
pretend all individuals are equally likely to live or die, *independent
of their species identity*, then you should expect this pattern
[@hubbell2001]. Later, John Harte showed that the principle of maximum
entropy was a more general approach to this question and could be
extended into other realms and patters.

As an example, consider a community of five individuals and three
species. Those are our only constraints, $N=5$, $S=3$. How many
different possible distributions of species abundances can we get if we
assume maximum uncertainty, i.e., maximum entropy? Is there any way to
know? Because we have so few, we can count them.

Consider five individuals distributed among three species (rows) and
list all possible community states (columns).


::: {.cell}

```{.r .cell-code}
# by default, R fills matrices by columns.
states <- matrix(
  c(3,1,1, # column 1,
    1,3,1, # column 2, etc.
    1,1,3,
    2,2,1,
    2,1,2,
    1,2,2),
  nrow = 3 )
# label species
rownames(states) <- LETTERS[1:3]
states # each column is a possible configuration.
```

::: {.cell-output .cell-output-stdout}

```
  [,1] [,2] [,3] [,4] [,5] [,6]
A    3    1    1    2    2    1
B    1    3    1    2    1    2
C    1    1    3    1    2    2
```


:::
:::


A rank or species abundance distributions does care *which* species is
most abundant, just that different species exist and can have different
abundances. Maximum entropy would have us ask, "what is the average
abundance of the most abundant species?" Let's calculate that. We can
see that the most abundant species has either 3 individuals or two
individuals and each of these occurs three times. Therefore the average
abundance of the most abundant species is calculated as


::: {.cell}

```{.r .cell-code}
(3+3+3+2+2+2)/6
```

::: {.cell-output .cell-output-stdout}

```
[1] 2.5
```


:::
:::


or two and half individuals out of five individuals in the whole
community.

How about the least abundant species? The least abundant species always
has only one individual. Therefore, the second most abundant species
has, on average, 1.5 individuals.


::: {.cell .column-margin}

```{.r .cell-code}
d <- data.frame(x=1:3, y=c(2.5, 1.5, 1) )

ggplot(d, aes(x, y) ) +
 geom_line() +  geom_point() + 
  labs(y="Expected abundance") 
```

::: {.cell-output-display}
![Expected rank abundance distribution, given five individuals and three species, under the maximum entropy theory of ecology.](01-theory_files/figure-pdf/fig-mete2-1.pdf){#fig-mete2 fig-pos='H'}
:::
:::


This is one prediction of METE: given only the number of individuals and
the number of species, it can predict a species abundance distribution.

## *In fine*

Useful theory helps us think more clearly and sometimes more
imaginatively about how the world works. Good theory is usually
*efficient* theory, making the greatest number of predictions with the
fewest free parameters. It helps clarify our thinking and disentangle
the myriad mechanisms at play in our ecological systems.


`<!-- quarto-file-metadata: eyJyZXNvdXJjZURpciI6Ii4ifQ== -->`{=html}

```{=html}
<!-- quarto-file-metadata: eyJyZXNvdXJjZURpciI6Ii4iLCJib29rSXRlbVR5cGUiOiJjaGFwdGVyIiwiYm9va0l0ZW1OdW1iZXIiOjMsImJvb2tJdGVtRmlsZSI6InJlZmVyZW5jZXMucW1kIiwiYm9va0l0ZW1EZXB0aCI6MH0= -->
```

# References 

```````{.quarto-title-block template='/Applications/quarto/share/projects/book/pandoc/title-block.md'}
---
title: References

---
\end{verbatim}

\phantomsection\label{refs}


\backmatter


\end{document}
